problem:  
Let \((M^n, g)\) be a complete, noncompact Riemannian manifold with **nonnegative sectional curvature**, and let \(S\subset M\) denote its **soul** in the sense of Cheeger–Gromoll.  
Suppose moreover that \(M\) has **Euclidean volume growth**,  
\[
\lim_{r\to\infty}\frac{\operatorname{Vol}(B(p,r))}{r^n}=v_0>0,
\]
and that its curvature tensor satisfies the **borderline pointwise decay**
\[
|Rm|(x)\le C(1+r(x))^{-2},\qquad r(x)=d(p,x),
\]
for some constant \(C>0\).

**Open Problem:**  
Does the borderline curvature decay \(r^{-2}\) combined with Euclidean volume growth constrain the intrinsic geometry or holonomy of the **normal bundle of the soul**?  

More precisely:
1. Must the soul \(S\) be **flat** under these hypotheses, or can it have positive curvature?  
2. If \(S\) is not flat, does this force the Gromov–Hausdorff tangent cones at infinity of \(M\) to be **nonunique** or **non-Euclidean**?  
3. Can one relate the **second fundamental form of \(S\)** in \(M\) or the curvature of the normal connection to the rate of curvature decay at infinity?

Equivalently, the problem seeks to determine whether **asymptotically flat geometry at infinity** propagates inward along the normal bundle to impose rigidity or flatness properties on the compact core.

Why is it a "good" problem:  
This version focuses on a subtle and unexplored interface between **asymptotic curvature decay** and the **intrinsic geometry of the soul**, which is untouched by the classical Soul Theorem.  
While the Soul Theorem gives a purely topological and qualitative description (\(M\cong\nu(S)\)), it provides no control on the metric geometry of \(S\) or its normal bundle. The current question asks whether quantitative decay of curvature at infinity forces **geometric rigidity of the compact core**—whether the “flatness at infinity” leaks back into the soul itself.

The problem ties together several deep areas:
- Large-scale geometry and asymptotic analysis (\(|Rm|\sim r^{-2}\) and Euclidean growth),
- Local differential geometry of totally geodesic submanifolds and their normal bundles,
- Metric limit and tangent-cone theory (possible nonuniqueness or shape constraints),
bridging the asymptotic and intrinsic scales of curvature structure.

The formulation is simple (a single decay assumption plus known topological facts) yet the implications reach to the heart of how asymptotic geometry controls global structure, an issue not resolved by current comparison or splitting results. Its resolution would reveal whether asymptotic flatness and curvature decay dictate the curvature and holonomy of the soul—an elegant, concise, and deeply open problem.